\section{Results and discussion}
\label{sec:results}

\subsection{Why threshold pruning fails}

Table~\ref{tab:v3diagnostic} records the observed v3 validation behavior. The
mean retained count is only moderately above the truth, but its distribution is
pathological: 49 of 128 curves retain all eight candidates and 22 retain none.
The result confirms that a mean-count loss does not constrain per-curve
cardinality and that globally coupled gate scores are poorly calibrated for a
fixed threshold.

\begin{table}[t]
  \centering
  \caption{Diagnostic results for the v3 threshold-pruning baseline on 128
  validation curves. These are not results of the proposed method.}
  \label{tab:v3diagnostic}
  \begin{tabular}{lr}
    \toprule
    Metric & Value \\
    \midrule
    True internal-knot count, mean & 3.281 \\
    Retained count at threshold 0.5, mean & 4.242 \\
    Curves retaining all 8 candidates & 49/128 \\
    Curves retaining zero candidates & 22/128 \\
    Knot precision / recall / F1 at 0.05 & 0.407 / 0.526 / 0.459 \\
    Matched-knot MAE & 0.0240 \\
    Standard B-spline refit RMS & 0.0160 \\
    \bottomrule
  \end{tabular}
\end{table}

\subsection{Main comparison}

The v4 categorical model reached a knot precision of approximately 0.44, but
its training and validation count accuracies were approximately 0.64 and 0.21,
respectively. Supplying the true count increased knot precision to approximately
0.55, whereas additionally supplying true point parameters produced no further
improvement. These diagnostics motivate canonical labels, ordinal local count
attention, and parameter sharing across count-conditioned representations.

\begin{table*}[t]
  \centering
  \caption{Main test-set comparison. Best values should be bold and the second
  best underlined only after all runs are complete.}
  \label{tab:main}
  \begin{tabular}{lrrrrrr}
    \toprule
    Method & Count acc. $\uparrow$ & Count MAE $\downarrow$ & Knot F1 $\uparrow$
    & Knot MAE $\downarrow$ & RMS $\downarrow$ & Time (ms) $\downarrow$ \\
    \midrule
    Fixed-$\Kmax$ & -- & \TODO{x} & \TODO{x} & \TODO{x} & \TODO{x} & \TODO{x} \\
    Independent gates & \TODO{x} & \TODO{x} & \TODO{x} & \TODO{x} & \TODO{x} & \TODO{x} \\
    Hard-Concrete v3 & \TODO{x} & \TODO{x} & \TODO{x} & \TODO{x} & \TODO{x} & \TODO{x} \\
    Count + top-$K$ & \TODO{x} & \TODO{x} & \TODO{x} & \TODO{x} & \TODO{x} & \TODO{x} \\
    Categorical count v4 & \TODO{x} & \TODO{x} & \TODO{x} & \TODO{x} & \TODO{x} & \TODO{x} \\
    Proposed ordinal, argmax & \TODO{x} & \TODO{x} & \TODO{x} & \TODO{x} & \TODO{x} & \TODO{x} \\
    Proposed ordinal, BIC & \TODO{x} & \TODO{x} & \TODO{x} & \TODO{x} & \TODO{x} & \TODO{x} \\
    \bottomrule
  \end{tabular}
\end{table*}

\TODO{Discuss whether improvements come from count accuracy, better knot
localization, or both. Include confidence intervals and statistical tests where
appropriate. Do not infer structural superiority from fitting error alone.}

\subsection{Robustness, generalization, and failure cases}

Evaluate increasing noise, nonuniform sampling, knot clustering, unseen control
point counts, and 3D curves. Show representative successes and failures,
including curves whose count is misclassified by one or more knots. Discuss
rank deficiency and the effect of control-polygon regularization.

\begin{figure}[t]
  \centering
  \fbox{\parbox[c][42mm][c]{0.90\linewidth}{\centering
  \TODO{Qualitative comparison: samples, true curve/knots, v3 threshold result,
  proposed result, and pointwise error plot.}}}
  \caption{Representative structural and geometric fitting results.}
  \label{fig:qualitative}
\end{figure}

\subsection{Limitations}

The supervised count formulation requires a specified geometric tolerance and
labeled knot vectors under a consistent parameterization. Greedy removal is
deterministic but is not guaranteed to find the globally smallest knot set.
The present study focuses on single open cubic
curves and does not yet handle closed curves, repeated knots, rational weights,
surface patch topology, or continuity constraints between CAD entities. A
an ordinal count error may still select the wrong representation, while BIC
selection increases deployment cost by refitting several complete branches.
